\chapter{Cryptography — RSA}
%%% generated by the blueprint pipeline from TCSlib.Cryptography.RSA
%%%%%%%%%%%%%%%%%%%%%%%%%%%%%%%%%%%%%%%%%%%%%%%%%%%%%%%%%%%%%%%%%%%%%%%%%%%%%%%
\section{Overview}

This file develops a small cost-tracking monad \texttt{TimeM T α}, whose values carry both a
return value of type $\alpha$ and an accumulated time cost of type $T$ (typically $T = \bbn$
counting comparisons), together with worked examples (median of three, insertion sort) that
prove functional correctness on the returned value and complexity bounds on the recorded cost.
It then formalizes the RSA cryptosystem: for distinct primes $p,q$ and exponents $e,d$ with
$ed = 1 + k(p-1)(q-1)$, decryption inverts encryption modulo $n = pq$.

%%%%%%%%%%%%%%%%%%%%%%%%%%%%%%%%%%%%%%%%%%%%%%%%%%%%%%%%%%%%%%%%%%%%%%%%%%%%%%%
\section{Declarations}

\begin{definition}[Pure computation in the time monad]
\lean{TimeM.pure}
\label{TimeM.pure}
\leanok
Lifts a value $a : \alpha$ into the computation $\langle a, 0\rangle$ of type \texttt{TimeM T α},
i.e.\ the computation that returns $a$ at zero time cost.
\end{definition}

\begin{definition}[Sequential composition in the time monad]
\lean{TimeM.bind}
\label{TimeM.bind}
\leanok
Given $m : \texttt{TimeM T α}$ and $f : \alpha \to \texttt{TimeM T β}$, the bind of $m$ and $f$
runs $f$ on the value returned by $m$ and returns $\langle (f\,m.\mathrm{ret}).\mathrm{ret},\;
m.\mathrm{time} + (f\,m.\mathrm{ret}).\mathrm{time}\rangle$: the two time costs are added.
\end{definition}

\begin{theorem}[Return value of a pure timed computation]
\lean{TimeM.ret_pure}
\label{TimeM.ret_pure}
\leanok
\difficulty{1}
A *timed computation* $\mathrm{TimeM}\,T\,\alpha$ records a return value in $\alpha$
together with an accumulated time cost in $T$. Suppose $T$ carries a distinguished zero
element, and let $\mathrm{pure}\,a$ denote the timed computation that returns $a$ at
zero cost. Then for every $a \in \alpha$, the return value of $\mathrm{pure}\,a$ is $a$.
\end{theorem}

\begin{theorem}[Return value of a bound cost-annotated computation]
\lean{TimeM.ret_bind}
\label{TimeM.ret_bind}
\leanok
\difficulty{1}
Consider computations that carry both a return value and an accumulated time cost drawn
from a type $T$ equipped with an addition, and let the monadic bind of such a
computation $m$ with a continuation $f$ run $m$, pass its return value to $f$, and add
the two costs. Then the return value of $m \mathbin{>\!\!>\!\!=} f$ equals the return
value of $f$ applied to the return value of $m$.
\end{theorem}

\begin{theorem}[Return value of a mapped timed computation]
\lean{TimeM.ret_map}
\label{TimeM.ret_map}
\leanok
\difficulty{1}
A value of type $\mathrm{TimeM}\ T\ \alpha$ is a timed computation, consisting of a
return value $\mathrm{ret} \in \alpha$ together with an accumulated time cost in $T$,
and $\mathrm{TimeM}\ T$ carries the functor structure that acts on the return value. For
any function $f : \alpha \to \beta$ and any timed computation $x$ of type
$\mathrm{TimeM}\ T\ \alpha$, the return value of the mapped computation obtained by
applying $f$ through this functor to $x$ equals $f(\mathrm{ret}(x))$.
\end{theorem}

\begin{theorem}[Return value of right-sequenced timed computations]
\lean{TimeM.ret_seqRight}
\label{TimeM.ret_seqRight}
\leanok
\difficulty{1}
Model a computation that tracks its running time as a pair consisting of a return value
and an accumulated time cost drawn from a type $T$ equipped with an addition. Given such
a computation $x$, with return value in a type $\alpha$, and a thunk $y$ producing a
second such computation with return value in a type $\beta$, the right-sequencing of $x$
and $y$ — which runs $x$, then $y()$, combines their time costs, and discards the value
of $x$ — has return value equal to that of $y()$.
\end{theorem}

\begin{theorem}[Return value of left-sequencing]
\lean{TimeM.ret_seqLeft}
\label{TimeM.ret_seqLeft}
\leanok
\difficulty{1}
A value of type $\mathrm{TimeM}\,T\,\alpha$ is a computation carrying both a return
value in $\alpha$ and an accumulated cost in $T$; write $\mathrm{ret}$ for the
return-value component. Suppose $T$ is equipped with an addition, let $x$ be such a
computation of type $\mathrm{TimeM}\,T\,\alpha$, and let $y$ be a (thunked) computation
of type $\mathrm{TimeM}\,T\,\beta$. Then the left-sequencing $x \mathbin{<{*}} y$, which
runs both computations and adds their costs but keeps only the value produced by $x$,
has return value $\mathrm{ret}(x \mathbin{<{*}} y) = \mathrm{ret}(x)$.
\end{theorem}

\begin{theorem}[Return value of sequenced timed computations]
\lean{TimeM.ret_seq}
\label{TimeM.ret_seq}
\leanok
\difficulty{1}
Regard a value of type $\mathrm{TimeM}\,T\,\alpha$ as a timed computation, carrying both
a return value in $\alpha$ and an accumulated time cost in $T$, and suppose $T$ is
equipped with an addition, so that such computations may be sequenced. Let $f$ be a
timed computation whose return value is a function $\alpha \to \beta$, and let $x$ be a
thunk, that is, a function on the one-point type producing a timed computation with
return value in $\alpha$. Then the return value of the computation obtained by
sequencing $f$ with $x$ equals the return value of $f$ applied to the return value of
the computation $x(\,)$ obtained by forcing the thunk.
\end{theorem}

\begin{theorem}[Additivity of time cost under bind]
\lean{TimeM.time_bind}
\label{TimeM.time_bind}
\leanok
\difficulty{1}
Model a computation as a pair consisting of a return value in a type $\alpha$ together
with an accumulated time cost drawn from a type $T$ equipped with an addition operation;
write $\mathrm{ret}(\cdot)$ for the returned value and $\mathrm{time}(\cdot)$ for the
cost. Given such a computation $m$ and a continuation $f$ that assigns to each value of
type $\alpha$ a further computation returning a value of type $\beta$, the sequential
composition $m \gg\!= f$ (run $m$, then run $f$ on its result) has cost equal to the
cost of $m$ plus the cost of running $f$ on the value returned by $m$: \[
\mathrm{time}(m \gg\!= f) = \mathrm{time}(m) +
\mathrm{time}\bigl(f(\mathrm{ret}(m))\bigr). \]
\end{theorem}

\begin{theorem}[Time cost of the pure computation]
\lean{TimeM.time_pure}
\label{TimeM.time_pure}
\leanok
\difficulty{1}
A cost-annotated computation over a type $T$ of costs and a value type $\alpha$ consists
of a return value in $\alpha$ together with an accumulated time cost in $T$. Suppose $T$
carries a distinguished zero element $0$. Then for every value $a \in \alpha$, the
trivial computation that returns $a$ (the $\mathrm{pure}$ computation) has time cost
$0$.
\end{theorem}

\begin{theorem}[Mapping preserves time cost]
\lean{TimeM.time_map}
\label{TimeM.time_map}
\leanok
\difficulty{1}
A timed computation over a cost type $T$ is a pair consisting of a return value in
$\alpha$ and an accumulated time cost in $T$. For any function $f\colon \alpha \to
\beta$ and any timed computation $x$, applying $f$ to the return value of $x$ (leaving
the value's type as $\beta$) yields a timed computation whose time cost equals that of
$x$.
\end{theorem}

\begin{theorem}[Time cost of right-sequencing]
\lean{TimeM.time_seqRight}
\label{TimeM.time_seqRight}
\leanok
\difficulty{1}
A *timed computation* over an additive cost type $T$ records a return value together
with an accumulated cost in $T$; write $\mathrm{time}(c) \in T$ for the cost of such a
computation $c$. Let $x$ be a timed computation and let $y$ be a thunk that, when
forced, produces a timed computation. The right-sequencing $x \mathbin{*>} y$, which
performs $x$ and then $y()$ and returns the value of $y()$, has time cost
\[\mathrm{time}(x \mathbin{*>} y) = \mathrm{time}(x) + \mathrm{time}(y()).\]
\end{theorem}

\begin{theorem}[Time cost of left-sequencing]
\lean{TimeM.time_seqLeft}
\label{TimeM.time_seqLeft}
\leanok
\difficulty{1}
Fix a type $T$ carrying an addition, and regard a value of type
$\mathrm{TimeM}\,T\,\alpha$ as a computation that records both a return value in
$\alpha$ and an accumulated time cost $\mathrm{time}\in T$. Given such a computation $x$
and a thunked computation $y\colon \mathrm{Unit}\to \mathrm{TimeM}\,T\,\beta$, the
left-sequencing $x \mathbin{<*} y$ — which runs both but keeps the return value of $x$ —
has time cost $\mathrm{time}(x) + \mathrm{time}\bigl(y()\bigr)$, the sum of the two
individual costs.
\end{theorem}

\begin{theorem}[Additivity of time cost under sequencing]
\lean{TimeM.time_seq}
\label{TimeM.time_seq}
\leanok
\difficulty{1}
Let $T$ be a type equipped with an addition, and regard a value of
$\mathrm{TimeM}\,T\,\alpha$ as a computation that returns a result in $\alpha$ together
with an accumulated time cost $\mathrm{time}(\cdot) \in T$. Given a computation $f$ of a
function $\alpha \to \beta$ and a thunked computation $x$ of an argument in $\alpha$,
the applicative sequencing $f \mathbin{\langle *\rangle} x$, which applies the computed
function to the computed argument, has time cost $\mathrm{time}(f) +
\mathrm{time}(x())$, the sum of the cost of $f$ and the cost of forcing and running $x$.
\end{theorem}

\begin{definition}[Tick]
\lean{TimeM.tick}
\label{TimeM.tick}
\leanok
The computation $\mathrm{tick}\,c$ returns the trivial value and records time cost $c$; it is the
basic way of charging $c$ units of time.
\end{definition}

\begin{theorem}[Return value of a tick]
\lean{TimeM.ret_tick}
\label{TimeM.ret_tick}
\leanok
\uses{TimeM.tick}
\difficulty{1}
For any time cost $c$, the return value of the computation $\mathrm{tick}\,c$ is the
trivial (unit) value.
\end{theorem}

\begin{theorem}[Time cost of a tick]
\lean{TimeM.time_tick}
\label{TimeM.time_tick}
\leanok
\uses{TimeM.tick}
\difficulty{1}
For any time cost $c$, the operation $\mathrm{tick}\,c$ — which charges $c$ units of
time and returns the trivial value — records accumulated time exactly $c$.
\end{theorem}

\begin{definition}[Maximum of two naturals, timed]
\lean{max2}
\label{max2}
\leanok
The timed computation returning the larger of two natural numbers $x,y$, charging one unit of
time for the single comparison $x \le y$.
\end{definition}

\begin{definition}[Median of three, timed]
\lean{median3}
\label{median3}
\leanok
For $a,b,c$ in a linear order, the timed computation returning a median of $a,b,c$ by a nested
sequence of comparisons, charging one unit of time per comparison performed.
\end{definition}

\begin{theorem}[Correctness of the timed median of three]
\lean{median3_correct}
\label{median3_correct}
\leanok
\uses{median3}
\difficulty{4}
Let $a$, $b$, $c$ be elements of a linearly ordered set. Then the value returned by the
timed computation $\mathrm{median3}(a,b,c)$ is a median of $a$, $b$, and $c$: it is one
of the three inputs occupying the middle position in the order, so that at least one of
the remaining two values is $\le$ it and at least one is $\ge$ it.
\end{theorem}

\begin{theorem}[Comparison count of the timed median of three]
\lean{median3_time}
\label{median3_time}
\leanok
\uses{median3}
\difficulty{4}
Let $a$, $b$, $c$ be elements of a linear order, and let $\mathrm{median3}(a,b,c)$ be
the timed computation that returns a median of the three by a nested sequence of
comparisons, charging one unit of time for each comparison performed. Then the recorded
time cost of $\mathrm{median3}(a,b,c)$ is at most $3$; that is, the computation performs
at most three comparisons.
\end{theorem}

\begin{definition}[Timed ordered insertion]
\lean{insert}
\label{insert}
\leanok
Inserts $x$ into a list, walking the list until the first element $y$ with $x \le y$ and charging
one unit of time per comparison. On a sorted input it produces the sorted list containing $x$.
\end{definition}

\begin{definition}[Timed insertion sort]
\lean{insertionSort}
\label{insertionSort}
\leanok
\uses{insert}
Insertion sort in the time monad: recursively sort the tail and then insert the head using the
timed insertion, accumulating the comparisons performed by every insertion step.
\end{definition}

\begin{definition}[Sortedness by adjacent pairs]
\lean{IsSorted}
\label{IsSorted}
\leanok
The predicate on lists defined recursively: the empty list and singletons are sorted, and
$x :: y :: \mathit{rest}$ is sorted when $x \le y$ and $y :: \mathit{rest}$ is sorted.
\end{definition}

\begin{theorem}[Adjacent sortedness coincides with being a $\le$-chain]
\lean{isSorted_iff_isChain}
\label{isSorted_iff_isChain}
\leanok
\uses{IsSorted}
\difficulty{4}
Let $\alpha$ be a linearly ordered type and let $\mathit{xs}$ be a finite list of
elements of $\alpha$. Then $\mathit{xs}$ is sorted — in the sense that each element is
at most the one following it — if and only if $\mathit{xs}$ is a chain for the relation
$\le$, meaning that $x \le y$ holds for every pair of consecutive entries $x, y$ of
$\mathit{xs}$.
\end{theorem}

\begin{theorem}[Adjacent sortedness equals pairwise sortedness]
\lean{isSorted_iff_sorted}
\label{isSorted_iff_sorted}
\leanok
\uses{IsSorted, isSorted_iff_isChain}
\difficulty{1}
Let $\mathit{xs}$ be a finite list over a linearly ordered set. Then $\mathit{xs}$ is
sorted in the adjacent-pairs sense if and only if it is sorted in the usual pairwise
sense with respect to $\le$; that is, if and only if $x_i \le x_j$ for every pair of
indices $i < j$.
\end{theorem}

\begin{theorem}[Correctness of timed ordered insertion]
\lean{ret_insert}
\label{ret_insert}
\leanok
\uses{insert}
\difficulty{3}
Let $\alpha$ be a linearly ordered type. For every $x \in \alpha$ and every finite list
$\mathit{xs}$ over $\alpha$, the value returned by the timed insertion of $x$ into
$\mathit{xs}$ (discarding the accumulated time cost) coincides with the ordinary ordered
insertion of $x$ into $\mathit{xs}$ with respect to the order $\le$: the list obtained
by placing $x$ immediately before the first entry $y$ of $\mathit{xs}$ satisfying $x \le
y$, and at the end if no such entry exists.
\end{theorem}

\begin{theorem}[Timed insertion sort returns the sorted list]
\lean{ret_insertionSort}
\label{ret_insertionSort}
\leanok
\uses{insertionSort, ret_insert}
\difficulty{3}
Let $\alpha$ be a linearly ordered type and let $xs$ be a list of elements of $\alpha$.
The value returned by the timed insertion sort of $xs$ equals the ordinary insertion
sort of $xs$ taken with respect to the order relation $\le$.
\end{theorem}

\begin{theorem}[Length of a timed insertion's output]
\lean{length_ret_insert}
\label{length_ret_insert}
\leanok
\uses{insert, ret_insert}
\difficulty{1}
Let $\alpha$ be a linearly ordered type, let $x \in \alpha$, and let $xs$ be a finite
list of elements of $\alpha$. The return value of the timed ordered insertion of $x$
into $xs$ is a list of length $\abs{xs} + 1$.
\end{theorem}

\begin{theorem}[Insertion sort preserves length]
\lean{length_ret_insertionSort}
\label{length_ret_insertionSort}
\leanok
\uses{insertionSort, ret_insertionSort}
\difficulty{1}
For every list $\mathit{xs}$ over a linearly ordered type, the output list produced by
the timed insertion sort of $\mathit{xs}$ has the same length as $\mathit{xs}$.
\end{theorem}

\begin{theorem}[Timed insertion returns a permutation]
\lean{insert_perm}
\label{insert_perm}
\leanok
\uses{insert, ret_insert}
\difficulty{1}
Let $\alpha$ be a linearly ordered type, let $x \in \alpha$, and let $xs$ be a finite
list of elements of $\alpha$. Then the list obtained by timed ordered insertion of $x$
into $xs$ is a permutation of $x :: xs$, the list with $x$ prepended to $xs$.
\end{theorem}

\begin{theorem}[Timed ordered insertion preserves sortedness]
\lean{insert_sorted}
\label{insert_sorted}
\leanok
\uses{IsSorted, insert, isSorted_iff_sorted, ret_insert}
\difficulty{2}
Let $\alpha$ be a linearly ordered type, let $x \in \alpha$, and let $\mathit{xs}$ be a
sorted list over $\alpha$. Then the list returned by the timed ordered insertion of $x$
into $\mathit{xs}$ — the return value of that computation, discarding its time cost — is
again sorted.
\end{theorem}

\begin{theorem}[Insertion sort returns a permutation of its input]
\lean{insertionSort_perm}
\label{insertionSort_perm}
\leanok
\uses{insertionSort, ret_insertionSort}
\difficulty{1}
Let $\alpha$ be a linearly ordered type and let $\mathit{xs}$ be a list of elements of
$\alpha$. Then the list produced by insertion sort applied to $\mathit{xs}$ is a
permutation of $\mathit{xs}$.
\end{theorem}

\begin{theorem}[Insertion sort produces a sorted list]
\lean{insertionSort_sorted}
\label{insertionSort_sorted}
\leanok
\uses{IsSorted, insertionSort, isSorted_iff_sorted, ret_insertionSort}
\difficulty{2}
Let $\alpha$ be a linearly ordered type. For every list $\mathit{xs}$ of elements of
$\alpha$, the list returned by the timed insertion sort of $\mathit{xs}$ is sorted.
\end{theorem}

\begin{theorem}[Correctness of timed insertion sort]
\lean{insertionSort_correct}
\label{insertionSort_correct}
\leanok
\uses{IsSorted, insertionSort, insertionSort_perm, insertionSort_sorted}
\difficulty{1}
Let $\alpha$ be a linearly ordered type and let $\mathit{xs}$ be a finite list of
elements of $\alpha$. Then the list returned by the timed insertion sort applied to
$\mathit{xs}$ is sorted, and it is a permutation of $\mathit{xs}$.
\end{theorem}

\begin{definition}[Insertion sort recurrence]
\lean{timeInsertionSortRec}
\label{timeInsertionSortRec}
\leanok
The worst-case comparison recurrence $T(0) = 0$ and $T(n+1) = T(n) + n$.
\end{definition}

\begin{theorem}[Cost of one ordered insertion]
\lean{time_insert_le}
\label{time_insert_le}
\leanok
\uses{insert}
\difficulty{3}
Let $\alpha$ carry a linear order, and for $x \in \alpha$ and a list $\mathit{xs}$ over
$\alpha$ let $\mathrm{insert}(x, \mathit{xs})$ denote the ordered insertion that scans
$\mathit{xs}$ from the front, spending one unit of time on each element compared, until
it reaches the first entry not smaller than $x$. Then the time cost incurred by
$\mathrm{insert}(x, \mathit{xs})$ is at most the length of $\mathit{xs}$.
\end{theorem}

\begin{theorem}[Insertion sort respects its comparison recurrence]
\lean{time_insertionSort_le_rec}
\label{time_insertionSort_le_rec}
\leanok
\uses{insert, insertionSort, ret_insertionSort, timeInsertionSortRec, time_insert_le}
\difficulty{4}
Let $\alpha$ be a linearly ordered type and let $xs$ be a finite list of elements of
$\alpha$. Running the timed insertion sort on $xs$ charges at most $T(\abs{xs})$ units
of time, where $\abs{xs}$ is the length of $xs$ and $T$ is the recurrence given by $T(0)
= 0$ and $T(n+1) = T(n) + n$.
\end{theorem}

\begin{theorem}[Quadratic bound for the insertion-sort recurrence]
\lean{timeInsertionSortRec_le_sq}
\label{timeInsertionSortRec_le_sq}
\leanok
\uses{timeInsertionSortRec}
\difficulty{3}
Let $T \colon \bbn \to \bbn$ be defined by the recurrence $T(0) = 0$ and $T(n+1) = T(n)
+ n$. Then for every $n \in \bbn$,
\[
T(n) \le n^2.
\]
\end{theorem}

\begin{theorem}[Quadratic time bound for insertion sort]
\lean{time_insertionSort_le_sq}
\label{time_insertionSort_le_sq}
\leanok
\uses{insertionSort, timeInsertionSortRec_le_sq, time_insertionSort_le_rec}
\difficulty{1}
Let $\mathit{xs}$ be a list of length $n$ over a linearly ordered type. Then the total
time cost accumulated by the timed insertion sort of $\mathit{xs}$ is at most $n^2$.
\end{theorem}

\begin{lemma}[Core RSA exponent identity]
\lean{RSA.rsa_core}
\label{RSA.rsa_core}
\leanok
\difficulty{4}
Let $p$ be a prime and let $c \in \bbn$ be arbitrary. Then for every $x \in \bbz/p\bbz$,
\[
x^{\,1 + c(p-1)} = x .
\]
This is the form of Fermat's little theorem underlying RSA, and it holds for all $x$,
including $x = 0$.
\end{lemma}

\begin{lemma}[RSA exponent identity modulo a prime]
\lean{RSA.rsa_zmod_of_factor}
\label{RSA.rsa_zmod_of_factor}
\leanok
\uses{RSA.rsa_core}
\difficulty{2}
Let $p$ be a prime, and let $m$, $N$, and $c$ be natural numbers satisfying $N = 1 +
c(p-1)$. Then
\[
m^{N} \equiv m \pmod{p};
\]
that is, $m^{N}$ and $m$ have the same image in $\bbz/p\bbz$. (The congruence holds for
every $m$, including multiples of $p$.)
\end{lemma}

\begin{lemma}[RSA exponentiation identity modulo a prime]
\lean{RSA.rsa_zmod_p}
\label{RSA.rsa_zmod_p}
\leanok
\uses{RSA.rsa_zmod_of_factor}
\difficulty{3}
Let $p$ be a prime, and let $q, m, e, d, k$ be natural numbers satisfying $ed = 1 +
k(p-1)(q-1)$. Then $m^{ed} \equiv m \pmod{p}$; that is, the images of $m^{ed}$ and $m$
agree in $\bbz/p\bbz$.
\end{lemma}

\begin{lemma}[RSA identity modulo the prime $q$]
\lean{RSA.rsa_zmod_q}
\label{RSA.rsa_zmod_q}
\leanok
\uses{RSA.PublicKey, RSA.rsa_zmod_of_factor}
\difficulty{1}
Let $q$ be a prime, and let $e, d, k, m, p$ be natural numbers satisfying $ed = 1 +
k(p-1)(q-1)$. Then $m^{ed} \equiv m \pmod{q}$; that is, the images of $m^{ed}$ and $m$
agree in $\bbz/q\bbz$.
\end{lemma}

\begin{lemma}[Chinese Remainder Theorem for RSA correctness]
\lean{RSA.rsa_crt}
\label{RSA.rsa_crt}
\leanok
\uses{RSA.PublicKey}
\difficulty{3}
Let $p$, $q$, $m$, and $ed$ be natural numbers with $p$ and $q$ coprime, and suppose
that $m^{ed} \equiv m \pmod{p}$ and $m^{ed} \equiv m \pmod{q}$. Then $m^{ed} \equiv m
\pmod{pq}$.
\end{lemma}

\begin{lemma}[CRT combination of a fixed-point congruence]
\lean{RSA.rsa_crt_pow}
\label{RSA.rsa_crt_pow}
\leanok
\uses{RSA.PublicKey, RSA.rsa_crt}
\difficulty{3}
Let $p, q, m, ed \in \bbn$ with $p$ and $q$ coprime, and write $\bar m$ for the residue
class of $m$ in each relevant quotient ring. If $\bar m^{\,ed} = \bar m$ in $\bbz/p\bbz$
and $\bar m^{\,ed} = \bar m$ in $\bbz/q\bbz$, then $\bar m^{\,ed} = \bar m$ in
$\bbz/pq\bbz$.
\end{lemma}

\begin{definition}[RSA public key]
\lean{RSA.PublicKey}
\label{RSA.PublicKey}
\leanok
A public key consists of a modulus $n \in \bbn$ and a public exponent $e \in \bbn$.
\end{definition}

\begin{definition}[RSA secret key]
\lean{RSA.SecretKey}
\label{RSA.SecretKey}
\leanok
\uses{RSA.PublicKey}
A secret key consists of a public key together with naturals $p, q, d, k$ and proofs that $p$ and
$q$ are prime, that $p \ne q$, that the public modulus satisfies $n = pq$, and that the exponents
satisfy $e d = 1 + k(p-1)(q-1)$.
\end{definition}

\begin{definition}[RSA encryption]
\lean{RSA.encrypt}
\label{RSA.encrypt}
\leanok
\uses{RSA.PublicKey}
Encryption of a message $m \in \bbn$ under a public key is $\bar m^{\,e} \in \bbz/n\bbz$, using
only the public data.
\end{definition}

\begin{definition}[RSA decryption]
\lean{RSA.decrypt}
\label{RSA.decrypt}
\leanok
\uses{RSA.SecretKey}
Decryption of a ciphertext $c \in \bbz/n\bbz$ under a secret key is $c^{\,d}$, where $d$ is the
private exponent.
\end{definition}

\begin{theorem}[Correctness of RSA]
\lean{RSA.rsa_correctness}
\label{RSA.rsa_correctness}
\leanok
\uses{RSA.SecretKey, RSA.decrypt, RSA.encrypt, RSA.rsa_crt_pow, RSA.rsa_zmod_p, RSA.rsa_zmod_q}
\difficulty{4}
Let $(n,e)$ be an RSA public key belonging to a secret key that also supplies distinct
primes $p$ and $q$ with $n = pq$, a private exponent $d$, and a natural number $k$
satisfying $ed = 1 + k(p-1)(q-1)$. Then for every message $m \in \bbn$, decrypting the
encryption of $m$ returns the class of $m$ modulo $n$: in $\bbz/n\bbz$ we have
$\bigl(\bar m^{\,e}\bigr)^{d} = \bar m$.
\end{theorem}

%%% added by the blueprint pipeline
%%%%%%%%%%%%%%%%%%%%%%%%%%%%%%%%%%%%%%%%%%%%%%%%%%%%%%%%%%%%%%%%%%%%%%%%%%%%%%%
\section{Additional declarations}

\begin{definition}[Timed computation]
\lean{TimeM}
\label{TimeM}
\leanok
For a cost type $T$ (typically $\bbn$) and a value type $\alpha$, a timed computation is
a pair consisting of a return value $\mathrm{ret} \in \alpha$ together with an
accumulated time cost $\mathrm{time} \in T$. This is the carrier of the time-complexity
monad: correctness properties are stated about the return value, and complexity bounds
about the recorded cost.
\end{definition}
